\documentclass[10pt,a4paper,ngerman]{scrartcl}
\usepackage[utf8]{inputenc}
\usepackage[ngerman]{babel}
\usepackage[T1]{fontenc}
\usepackage{amsmath}
\usepackage{amsfonts}
\usepackage{amssymb}
\usepackage{makeidx}
\usepackage{graphicx}
\usepackage{epstopdf}
\usepackage[onehalfspacing]{setspace}
\usepackage{array}
\usepackage{upgreek}
\usepackage{float}
\usepackage{csquotes}
\usepackage{chemformula}
\usepackage{chemgreek}
\usepackage{chemmacros}
\chemsetup{modules=all}
\usepackage{tikz}
\usepackage{siunitx}
\sisetup{locale = DE,per-mode = symbol}
\usepackage{esvect}
\usepackage{geometry}
\usepackage{pdfpages}
\usepackage[font=small]{caption}
%\usepackage[version=4]{mhchem}
\usepackage[backend=biber, style=chem-angew]{biblatex}
\addbibresource{DSC.bib}
\usepackage{tabularx, booktabs, multirow}
\usepackage[pdfborder={0 0 0}]{hyperref}
\usepackage{scrlayer-scrpage}%Header für die KOMA-script -Klasse
%\usepackage{hyperref}
\usepackage{pgfplots}
\usepackage{textgreek}


\newcommand{\tildenu}{{\fontencoding{LGR}\selectfont\accperispomeni\textnu}}
\captionsetup{format=plain}
\chemsetup[phases]{pos=sub}
\chemsetup[redox]{pos=top}
\chemsetup[redox]{align=center}
%\chemsetup[redox]{explicit-sign=true}
%\chemsetup[redox]{roman=false}
\newcommand{\celsius}{^{\circ}\mathrm{C}}
\parindent0pt
\sloppy
\DeclareChemPhase{\s}{s}
\DeclareChemPhase{\l}{l}
\DeclareChemPhase{\g}{g}
\renewcaptionname{ngerman}{\figurename}{Abb.}
\renewcaptionname{ngerman}{\tablename}{Tab.}
\geometry{bottom=100pt} \geometry{top=80pt}
\newcommand{\m}{\mathrm{m}}
\newcommand{\K}{\mathrm{K}}
\newcommand{\J}{\mathrm{J}}
\newcommand{\gr}{\mathrm{g}}
\newcommand{\kg}{\mathrm{kg}}
\newcommand{\mol}{\mathrm{mol}}
\newcommand{\Pa}{\mathrm{Pa}}
\newcommand{\ad}{\mathrm{ad}}
\newcommand{\de}{\mathrm{de}}
\newcommand{\gmol}{\mathrm{\frac{g}{mol}}}
\newcommand{\JKmol}{\mathrm{\frac{J}{K \cdot mol}}}
\newcommand{\kgmss}{\mathrm{\frac{kg}{m \cdot s^2}}}
\newcommand{\loge}{\mathrm{ln}}



\begin{document}
\begin{titlepage}
	\vspace*{2,5cm}
	\begin{center}
		\begin{LARGE}
			\textbf{Polymer Chemie}\\
			\textbf{Polyaddition und Polykondensation}\\
		\end{LARGE}
		\vspace{1cm}
		\textbf{Universität Stuttgart}\\
		\vspace*{1,5cm}
		\begin{tabular}{lp{7,5cm}l}
			Author 1:     &Laurens Viehoff\\
            Author 2:     &Felix Roth\\   
			E-Mail 1:      &st183688@stud.uni-stuttgart.de\\
            E-Mail 2:      &st188157@stud.uni-stuttgart.de  \\
			Student-ID 1: & 3676761   \\ 
            Student-ID 2: & 3711192 \\ \\
			Assistentin:  &   \\ \\ 
		\end{tabular}
	\end{center}

    \begin{center}
        \textbf{\Large Abstract}
    \end{center}
    
\end{titlepage}
\thispagestyle{empty}
\renewcommand{\thepage}{\arabic{page}}
\newpage
\tableofcontents
\setcounter{page}{1}
\renewcommand{\thepage}{\arabic{page}}
\newpage

\section{Theorie}

    

\section{Durchführung}

	\subsection{Herstellung von Nylon}

		Es wurde eine Hexamethylendiamin Lösung (\SI{2,2}{g}, \SI{13}{mmol}, \SI{50}{mL} Wasser) hergestellt und in einem Becherglas vorgelegt.
		Anschließend wurde eine weitere Lösung Sebacinsäuredichlorid (\SI{1,5}{mL}, \SI{7}{mmol}, \SI{100}{mL} Cyclohexan) hergestellt.
		Die erste der beiden Lösungen wurde zusätzlich mit etwas Phenolphtalein angefärbt, sodass die Phasengrenze besser zu sehen ist.
		Daraufhin wurde die Hexamethylendiamin Lösung mit der zweiten Lösung vorsichtig überschichtet. 
		Nach überschichten wurde eine pinzette und ein glasstab verwended um einen faden aus der Phasengrenze zu ziehen.
		Falls das ziehen eines Fadens möglich ist wird dieser anschließend mit Natriumhydrogencarbonatlösung, Wasser und Aceton gewaschen und zum trocknen in den Trockenschrank gelegt.

	\subsection{Lösungspolymerisation}

		Es wird Polytetrahydrofuran (Mn ~ 2000, \SI{4}{g}, \SI{2}{mmol}, 1 Äq) in einem schlenkkolben vorgelegt, und anschließend ausgeheizt.
		Es wird TDI (\SI{0,8}{g}, \SI{0,66}{mL}, \SI{4,6}{mmol}, 2,3 Äq), \SI{40}{mL} trockenes DMF sowie eine DBTL (Dibutyzinnlaureat) Lösung (\SI{4}{mL}, \SI{0,1}{mol/L}, in DMF) 
		zugegeben und anschließend auf \SI{80}{°C} für eine stunde, unter rühren erhitzt.
		Danach wird 1,4-Butandiol (\SI{0,23}{g}, \SI{0,23}{mL}, \SI{2,6}{mmol}) zugegeben und weitere zwei stunden gerührt.
		Die lösung wird nach den zwei Stunden in MeOH (250mL) gegeben und im falle von ausfallendem Polymer wird dieses abgefiltert und mit MeOH gewaschen.
		Das Polymer wird im trochenschrank getrocknet.


\printbibliography

\end{document}